\chapter{Reference Indexing}
\label{ch:indexing}

The index (\code{SketchIndex}) is built exactly once per run, from the reference FASTA, before any
read is processed. Its job is to answer, in $O(1)$ average time, two questions for any $k$-mer
hash $h$: \emph{how many times does $h$ occur in the reference} (\S\ref{sec:index-count}), and
\emph{at which (segment, position) pairs} (used throughout bucketing/scoring).

\section{Data owned by the index}

\begin{center}
\begin{tabularx}{\textwidth}{@{}llX@{}}
\toprule
\textbf{Field} & \textbf{Type} & \textbf{Meaning} \\
\midrule
\code{segments} & \code{Vec<RefSegment>} & One entry per FASTA record, in file order
  (\S\ref{sec:type-refsegment}). \\
\code{h2single} & \code{hash\_t} $\to$ \code{Hit} & $k$-mers occurring \emph{exactly once} across
  the whole reference. \\
\code{h2multi} & \code{hash\_t} $\to$ \code{Vec<Hit>} & $k$-mers occurring \emph{two or more}
  times, each list sorted (\S\ref{sec:index-sort}). \\
\bottomrule
\end{tabularx}
\end{center}

Splitting single- vs.\ multi-occurrence $k$-mers into two separate maps (rather than one map to
always-a-list) is a memory/speed optimization: the overwhelming majority of $k$-mers in a
non-repetitive genome occur exactly once, so storing those as a bare \code{Hit} (no vector
indirection/allocation) rather than a one-element \code{Vec<Hit>} meaningfully reduces both memory
and pointer-chasing.

\section{Building the index: \code{build\_index}}

\begin{algobox}{title={\code{build\_index}(t\_file, k, h\_frac, max\_matches)}}
\begin{verbatim}
for (segm_name, seq) in read_fasta(t_file):
    sketch := FracMinHash.sketch(seq, k, h_frac)          # Chapter 5
    segm_id := len(segments)
    segments.append(RefSegment{ kmers: sketch, name: segm_name,
                                 sz: len(seq), id: segm_id })
    populate_h2pos(sketch, segm_id)                        # below

# migrate any hash that ended up in BOTH maps fully into h2multi
for (h, hits) in h2multi:
    if h in h2single:
        hits.append(h2single.remove(h))
    sort hits by (segm_id, r)                               # enables binary search later

if max_matches is set:
    erase_frequent_kmers(max_matches)                       # below
\end{verbatim}
\end{algobox}

\subsection{\code{populate\_h2pos}: routing each hit to the right map}
\label{sec:index-populate}

\begin{algobox}{title={\code{populate\_h2pos}(sketch, segm\_id)}}
\begin{verbatim}
for tpos, kmer in enumerate(sketch):
    hit := Hit(kmer, tpos, segm_id)
    if kmer.h not in h2single:
        h2single[kmer.h] := hit                 # first-ever occurrence: goes to h2single
    else:
        # second-or-later occurrence: goes to h2multi (h2single's own
        # earlier entry is migrated into h2multi at the *end* of indexing,
        # not immediately -- see build_index above)
        if max_matches is unset or len(h2multi[kmer.h]) < max_matches + 1:
            h2multi[kmer.h].append(hit)
\end{verbatim}
\end{algobox}

\begin{notebox}
Note the asymmetry this creates \emph{during} indexing (before the end-of-\code{build\_index}
migration step runs): a hash with exactly 2 occurrences temporarily has its \emph{first}
occurrence sitting in \code{h2single} and only its \emph{second} occurrence in \code{h2multi} ---
\code{h2multi[h]} would report length 1, not 2, if queried mid-build. This is why the migration
step (moving \code{h2single}'s entry into \code{h2multi} and \emph{erasing} it from
\code{h2single}) must run once, after all segments have been sketched, before the index is
considered complete and used for querying.
\end{notebox}

The \code{max\_matches + 1} cap during population (rather than exactly \code{max\_matches}) is
deliberate slack: it lets \code{erase\_frequent\_kmers} (below) later distinguish ``exactly at the
cap'' from ``over the cap'' after the single-entry migration adds one more; in practice, since
\code{erase\_frequent\_kmers} deletes any hash whose \emph{final} count exceeds \code{max\_matches}
outright, the effect is simply an upper bound on wasted work spent appending hits for a hash that
is going to be discarded anyway.

\subsection{Sorting multi-hit lists}
\label{sec:index-sort}

Each \code{h2multi} list is sorted by \code{(segm\_id, r)} lexicographically. Because $r$ (a
hit's nucleotide position) and $\mathit{tpos}$ (its index into the segment's own sketch array)
increase together within one segment (the sketch is built and stored in increasing position
order, \S\ref{sec:type-refsegment}), sorting by $r$ \emph{also} sorts by $\mathit{tpos}$
consistently --- so binary search by either coordinate space works correctly against the same
sorted list, which matters because bucket/window boundary checks are expressed in \emph{either}
space depending on the \code{abs\_pos} flag (\S\ref{sec:config-axes}).

\subsection{\code{erase\_frequent\_kmers}: capping pathologically repetitive $k$-mers}
\label{sec:max-matches}

CLI flag \code{-M} sets \code{max\_matches}, an optional hard cap on how many reference
occurrences a single $k$-mer hash is allowed to have before it is discarded from the index
entirely (treated thereafter as if it never occurred in the reference at all, i.e.\
\code{hits\_in\_t} $=0$ for it). This exists because extremely repetitive $k$-mers (satellite
repeats, low-complexity runs) contribute enormous numbers of hits that are almost never
informative for mapping but are expensive to iterate during seeding/bucketing --- capping them is
a standard technique shared with essentially every other $k$-mer-based mapper (minimap2's
\texttt{-f}, Winnowmap's repetitive-$k$-mer blacklist, etc.).

\begin{algobox}{title={\code{erase\_frequent\_kmers}(max\_matches)}}
\begin{verbatim}
for (h, hits) in h2multi:
    if len(hits) > max_matches:
        delete h2multi[h]
\end{verbatim}
\end{algobox}

Note this only prunes \code{h2multi} (by construction, \code{h2single} entries always have exactly
one occurrence, so they can never exceed any cap $\ge 1$).

\section{\code{count(h)}: reference occurrence count for a hash}
\label{sec:index-count}

\begin{algobox}{title={\code{count}(h)}}
\begin{verbatim}
if h in h2single: return 1
if h in h2multi:  return len(h2multi[h])
return 0
\end{verbatim}
\end{algobox}

This is the function seeding (Chapter~\ref{ch:seeding}) calls once per distinct read $k$-mer to
determine \code{hits\_in\_t} (\S\ref{sec:type-seed}) --- it is the quantity seeds are sorted by
(rarest first) and the input to the seed-heuristic bound's implicit assumption that
\code{hits\_in\_t}$=0$ seeds contribute nothing.

\section{Reference C++ implementation}

\begin{lstlisting}[style=cppstyle,caption={Core of \code{SketchIndex}, \texttt{src/index.h}}]
inline rpos_t count(hash_t h) const {
    if (auto it = h2single.find(h); it != h2single.end()) return 1;
    if (auto it = h2multi.find(h); it != h2multi.end()) return it->second.size();
    return 0;
}

void populate_h2pos(const sketch_t& sketch, segm_t segm_id) {
    for (size_t tpos = 0; tpos < sketch.size(); ++tpos) {
        const Kmer& kmer = sketch[tpos];
        const auto hit = Hit(kmer, tpos, segm_id);
        if (!h2single.contains(kmer.h))
            h2single[kmer.h] = hit;
        else if (H->params.max_matches == -1
                 || (rpos_t)h2multi[kmer.h].size() < H->params.max_matches + 1)
            h2multi[kmer.h].push_back(hit);
    }
}

void build_index(const std::string &tFile) {
    read_fasta_klib(tFile, [this](const string &segm_name, const string &T, float) {
        sketch_t t = H->sketcher.sketch(T);
        add_segment(segm_name, T, t);   // pushes RefSegment, calls populate_h2pos
    });

    for (auto &[h, hits] : h2multi) {
        if (h2single.contains(h)) {
            hits.push_back(h2single.at(h));
            h2single.erase(h);
        }
        sort(hits.begin(), hits.end(), [](const Hit &a, const Hit &b) {
            if (a.segm_id != b.segm_id) return a.segm_id < b.segm_id;
            return a.r < b.r;
        });
    }
    if (H->params.max_matches != -1) erase_frequent_kmers();
}
\end{lstlisting}

\section{Complexity and memory}

Indexing costs $O(n)$ time for sketching (Chapter~\ref{ch:sketching}) plus $O(rn \log d)$ for
sorting, where $rn$ is the total number of kept $k$-mers and $d$ is the maximum multiplicity of
any single hash (the sort is per-hash, so the $\log d$ factor applies independently to each
list, not globally). Memory is $O(rn)$: one map entry per distinct occurrence, plus $O(rn)$ for the
stored per-segment sketches themselves (each \code{RefSegment.kmers}).
